\section{范数}

% 在当前项目中使用的是 ctexart 文档类:
% Markdown 的 #   对应 \section
% Markdown 的 ##  对应 \subsection
% Markdown 的 ### 对应 \subsubsection
% 更细的层级通常先用段落文字或列表组织,必要时再用 \paragraph.

\subsection{基本定义}

设 $V$ 是数域 $\K$ 上的线性空间.若映射
\[
    \|\cdot\|:V\to\R
\]
满足以下三条性质,则称 $\|\cdot\|$ 为 $V$ 上的一个范数:
\begin{enumerate}
    \item 正定性: $\|x\|\ge 0$,且 $\|x\|=0$ 当且仅当 $x=0$;
    \item 齐次性: 对任意 $\alpha\in\K$ 和 $x\in V$,有 $\|\alpha x\|=|\alpha|\|x\|$;
    \item 三角不等式: 对任意 $x,y\in V$,有 $\|x+y\|\le \|x\|+\|y\|$.
\end{enumerate}

\subsection{常见例子}

\subsubsection{向量的 \texorpdfstring{$p$}{p}-范数}

对 $x=(x_1,\dots,x_n)^\top\in\K^n$,当 $1\le p<\infty$ 时定义
\[
    \|x\|_p=\left(\sum_{i=1}^n |x_i|^p\right)^{1/p}.
\]
特别地,$p=2$ 时得到欧氏范数:
\[
    \|x\|_2=\sqrt{x^\ast x}.
\]

\subsubsection{无穷范数}

向量 $x\in\K^n$ 的无穷范数定义为
\[
    \|x\|_\infty=\max_{1\le i\le n}|x_i|.
\]
它刻画的是向量各分量绝对值中的最大值.

\subsection{一个基本性质}

\begin{lemma}
若 $\|\cdot\|$ 是 $V$ 上的范数,则对任意 $x,y\in V$ 有
\[
    \bigl|\|x\|-\|y\|\bigr|\le \|x-y\|.
\]
\end{lemma}

\begin{proof}
由三角不等式,
\[
    \|x\|=\|(x-y)+y\|\le \|x-y\|+\|y\|,
\]
故 $\|x\|-\|y\|\le \|x-y\|$.交换 $x,y$ 可得
\[
    \|y\|-\|x\|\le \|y-x\|=\|x-y\|.
\]
合并两式即得结论.
\end{proof}
